\documentclass[10pt,twocolumn]{article}

\usepackage[a4paper,margin=1.8cm]{geometry}
\usepackage{graphicx}
\usepackage{amsmath}
\usepackage{amssymb}
\usepackage{booktabs}
\usepackage{hyperref}
\usepackage{float}
\usepackage{caption}

\title{\textbf{Practical Work 3: Segmentation of COVID-19 X-ray Images}}
\author{Tran Quang Huy - 23BI14196}
\date{January 2026}

\begin{document}

\maketitle

\section{Introduction}

Image segmentation has a crucial role in medical diagnosis. For example, during the COVID-19 pandemic, automatic detection and segmentation tools are essential for doctors to indentify and treat infected lung areas.

In this practical work, we attempted to segment infected regions of COVID-QU-Ex dataset \cite{covid_qu_ex} using the U-Net architecture. The implementation follows the tutorial by PyImageSearch \cite{pyimagesearch_unet}, which was also introduced in our Introduction to Deep Learning course.

\section{Dataset Description}

The COVID-QU-Ex dataset contains chest X-ray images of COVID-19 patients along with corresponding infection masks.

We resized all images to 256 x 256 and converted to grayscale. 
Pixel intensities were normalized to the range [0,1].

The number of training images are 1864, and 466 for validation images.

\section{Methodology}

\subsection{U-Net Architecture}

We implemented an U-Net architecture, more details can be found in the link below, and all credits go to Margaret Maynard-Reid. 

In a nutshell, segmentation model like U-Net learns to classify and predict each pixel, but in training each pixels won't be treated individually but together in order to get all the context from the image.

\subsection{Training Details}

    We used a batch size of 8 and a number epoch of 20. Training with a high number of epochs will demonstrate accurately the capability of the U-Net architecture.

\section{Experiments and Results}

\subsection{Training Performance}

After training for 20 epochs:

\begin{itemize}
    \item Final Training Accuracy: \textbf{92\%}
    \item Final Validation Accuracy: \textbf{92\%}
    \item Final Training Loss: \textbf{0.17}
    \item Final Validation Loss: \textbf{0.15}
\end{itemize}

Good results overall.

\subsection{Qualitative Results}

Figure~\ref{fig:results} shows example segmentation outputs.

\begin{figure}[H]
    \centering
    \includegraphics[width=0.9\linewidth]{Screenshot 2026-03-02 194405.png}
    \caption{Input image, real mask, and predicted mask.}
    \label{fig:results}
\end{figure}

The model successfully show that there are infection regions, although not entirely accurate.

\section{Comparison with State of the Art}

A study using an advanced U-Net variant\cite{optimizedUnet} net a results of: 

\begin{itemize}
    \item Mean Intersection over Union: \textbf{82.72\%}
    \item Mean Dice Coefficient: \textbf{78.37\%}
    \item Accuracy: \textbf{95.17\%}
\end{itemize}

Compared to this, our implementation
used a relatively simple U-Net. While it performed reasonable
well, it may not reach the performance of more advanced architectures.

\section{Conclusion}

In this practical work, we implemented a segmentation model 
for COVID-19 infection detection in X-ray images using simple U-Net architecture. 

The model successfully learned to identify infection regions and demonstrated the usefulness of automatic segmentation tools in medical imaging tasks.

\begin{thebibliography}{9}

\bibitem{covid_qu_ex}
COVID-QU-Ex, 
``consists of 33,920 chest X-ray (CXR) images'' 
Available at: \url{https://www.kaggle.com/datasets/anasmohammedtahir/covidqu}

\bibitem{pyimagesearch_unet}
PyImageSearch, 
``U-Net Image Segmentation in Keras,'' 
Available at: \url{https://pyimagesearch.com/2022/02/21/u-net-image-segmentation-in-keras/}

\bibitem{optimizedUnet}
Optimized attention U-Net for enhanced lung and area of infection segmentation in chest X-Rays and CT scans. 
Available at: \url{https://www.sciencedirect.com/science/article/pii/S1687850725003620?utm_source=chatgpt.com}

\end{thebibliography}

\end{document}